@ID: OW21D1P1I
@BIDANG: Analisis Real

Diberikan himpunan $A\subset\mathbb R$ dengan

$$
B=\left\{\frac1x:x\in A\right\}.
$$

Jika $\inf(A)>0$, maka $\sup(B)=\ldots$, dan jika $\inf(A)=0$, maka $\sup(B)=\ldots$

@ID: OW21D1P2I
@BIDANG: Analisis Real

Diberikan barisan bilangan real $(a_n)$ dengan

$$
a_n=\sum_{k=1}^{n}\frac{k!+4k+1}{k!\,4^k},
$$

maka

$$
\lim_{n\to\infty}a_n=\ldots
$$

@ID: OW21D1P3I
@BIDANG: Struktur Aljabar

Misalkan $a$ adalah bilangan asli terkecil yang mengakibatkan

$$
\mathbb Z[x]/\langle a,x^2+1\rangle
$$

merupakan suatu lapangan yang mempunyai lebih dari satu unsur. Jika banyaknya unsur pada lapangan ini adalah $b$, maka nilai $a\times b=\ldots$

@ID: OW21D1P4I
@BIDANG: Struktur Aljabar

Misalkan $G$ adalah suatu grup berorde $2021$. Misalkan juga $x$ dan $y$ merupakan dua unsur di $G$ yang tidak sama dengan identitas dan memiliki orde yang berbeda. Jika $H$ adalah subgrup terkecil yang memuat $x$ dan $y$, maka banyaknya unsur di $H$ adalah $\ldots$

@ID: OW21D1P5I
@BIDANG: Kombinatorika

Solusi dari relasi rekuren

$$
x_n=4x_{n-1}-3x_{n-2}+2^n,\qquad (n\ge3)
$$

dengan syarat $x_1=1$ dan $x_2=11$ adalah $\ldots$

@ID: OW21D1P1U
@BIDANG: Analisis Real

Diberikan barisan bilangan real positif $(x_n)$ dengan

$$
x_{n+1}\le x_n-\frac1{2^n},\qquad n\in\mathbb N.
$$

Buktikan bahwa barisan $(x_n)$ konvergen, kemudian tentukan nilai limitnya.

@ID: OW21D1P2U
@BIDANG: Analisis Real

Diberikan fungsi kontinu $f:\mathbb R\to\mathbb R$. Jika setiap himpunan terbuka $G$ diperoleh $f(G)$ terbuka, buktikan bahwa $f$ merupakan fungsi monoton.

@ID: OW21D1P3U
@BIDANG: Struktur Aljabar

Misalkan $S$ adalah suatu himpunan yang memiliki dua operasi biner $\circ$ dan $*$. Diketahui bahwa masing-masing operasi mempunyai unsur identitas (yang tidak mesti sama) dan untuk setiap $a,b,c,d\in S$ berlaku

$$
(a*b)\circ(c*d)=(a\circ c)*(b\circ d).
$$

Haruskah operasi biner $\circ$ dan $*$ merupakan operasi yang sama?

@ID: OW21D1P4U
@BIDANG: Struktur Aljabar

Misalkan $\mathbb Z_3[x]$ merupakan gelanggang polinom dengan koefisien di $\mathbb Z_3$. Diberikan $f(x)\in\mathbb Z_3[x]$ dengan $x^{2k}-1$ membagi $(f(x))^2-1$ untuk suatu bilangan asli $k\ge1$. Buktikan:

(a) $(f(x)-1)(f(x)+1)=a(x)(x^2-1)$ dan $(f(x)-x)(f(x)+x)=b(x)(x^2-1)$ untuk suatu $a(x),b(x)\in\mathbb Z_3[x]$.

(b) Jika $x^2-1$ tidak membagi $f(x)-x$ dan $f(x)+x$, maka $x^2-1$ membagi $f(x)-1$ atau $f(x)+1$.

@ID: OW21D1P5U
@BIDANG: Analisis Real

Misalkan $N$ adalah suatu bilangan bulat positif dan $\xi$ adalah suatu bilangan irasional. Buktikan bahwa terdapat bilangan rasional $\frac ab$ dengan $1\le b\le N$ yang memenuhi

$$
\left|\xi-\frac ab\right|
<
\frac1{b(N+1)}.
$$

@ID: OW21D2P1I
@BIDANG: Aljabar Linear

Misalkan $T$ pemetaan linear dari $\mathbb R^{2021}$ ke $\mathbb R^3$ sehingga untuk setiap $x\in\mathbb R^{2021}$, vektor $T(x)$ berbentuk

$$
T(x)=
\begin{pmatrix}
a\\
b\\
a+b
\end{pmatrix}
$$

untuk suatu $a,b\in\mathbb R$. Dimensi terkecil yang mungkin dari $\ker(T)$ adalah $\ldots$

@ID: OW21D2P2I
@BIDANG: Aljabar Linear

Banyaknya nilai eigen positif dari matriks

$$
\begin{pmatrix}
-1&3&0&0&0\\
3&2&-1&0&0\\
0&-1&2&-1&0\\
0&0&-1&2&-1\\
0&0&0&-1&-2
\end{pmatrix}
$$

adalah $\ldots$

@ID: OW21D2P3I
@BIDANG: Analisis Kompleks

Jika $a$ dan $b$ merupakan bilangan real dengan $0<a<b<2\pi$ yang memenuhi persamaan

$$
e^{ai}+e^{bi}=i\sqrt2,
$$

maka nilai $\dfrac ba=\ldots$

@ID: OW21D2P4I
@BIDANG: Analisis Kompleks

Fungsi kompleks $f(z)=\dfrac1z$ memetakan garis $\operatorname{Re}(z)=\dfrac14$ ke lingkaran dengan jari-jari $\ldots$

@ID: OW21D2P5I
@BIDANG: Kombinatorika

Kata sandi tanpa perulangan karakter dibentuk menggunakan huruf kapital. Sebuah kata sandi dikatakan sempurna bila tidak memuat untaian karakter **XYZ** maupun **ZYX**. Besarnya peluang untuk membentuk kata sandi sempurna yang terdiri atas $8$ huruf adalah $\ldots$

@ID: OW21D2P1U
@BIDANG: Aljabar Linear

(a) Misalkan $T:\mathbb R^3\to\mathbb R^3$ didefinisikan dengan

$$
T(x)=
\begin{cases}
-x,&\text{jika }x=(1,0,0),\\
x,&\text{jika }x\ne(1,0,0).
\end{cases}
$$

Buktikan bahwa $T$ bukan pemetaan linear dan jika $u\cdot v=0$, maka berlaku $T(u)\cdot T(v)=0$.

(b) Jelaskan apakah terdapat pemetaan tidak linear $T:\mathbb R^3\to\mathbb R^3$ sehingga berlaku

$$
T(u)\cdot T(v)=u\cdot v
$$

untuk setiap $u,v\in\mathbb R^3$.

@ID: OW21D2P2U
@BIDANG: Aljabar Linear

Suatu matriks $B$ berukuran $2\times2$ dengan komponen real dikatakan menarik jika terdapat matriks $A$ berukuran $2\times2$ dengan komponen real sehingga

$$
AB-BA=B^2.
$$

(a) Jelaskan apakah

$$
\begin{pmatrix}
0&1\\
0&0
\end{pmatrix}
$$

merupakan matriks menarik.

(b) Jika matriks $B$ menarik, buktikan bahwa $B$ nilpoten, yaitu $B^2=0$.

@ID: OW21D2P3U
@BIDANG: Analisis Kompleks

Misalkan

$$
A=\{z\in\mathbb C:|z|\le1\}
$$

dan

$$
B=\{\,x+(4-2x)i:x\in\mathbb R\,\}.
$$

Jika

$$
C=\{a-b:a\in A,\ b\in B\},
$$

tentukan

$$
\inf\{|c|:c\in C\}.
$$

@ID: OW21D2P4U
@BIDANG: Analisis Kompleks

Diketahui fungsi $f(z)$ analitik pada domain $D$. Jika ada konstanta $c_1,c_2\in\mathbb C$ yang tidak semuanya nol sehingga

$$
c_1f(z)+c_2\overline{f(z)}=0
$$

untuk setiap $z\in D$, maka buktikan bahwa $f(z)$ merupakan fungsi konstan pada $D$.

@ID: OW21D2P5U
@BIDANG: Kombinatorika

Pengubinan dengan panjang $n\ge1$ adalah sebuah cara menutupi lantai berukuran $1\times n$ menggunakan ubin $1\times1$ berwarna merah, putih, hijau, kuning, atau biru.

Sebuah pengubinan dikatakan ideal bila menggunakan genap ubin merah, genap ubin hijau, dan ganjil ubin biru.

Tentukan banyaknya pengubinan ideal dengan panjang $n$.